\chapter{Теоретические основы и проектирование мультиагентных систем навигации}
\label{chap:theory}

В данной главе рассматриваются теоретические аспекты функционирования мультиагентных систем навигации, проводится критический анализ современных подходов к моделированию транспортных потоков, а также проектируются математические и архитектурные решения для высокопроизводительной динамической маршрутизации реального времени.

\section{Сравнительный обзор подходов к мультиагентному моделированию и навигации}

При проектировании навигационных комплексов для городских агломераций сверхбольшой размерности ($\approx 600$ тыс. сегментов дорожной сети, $\approx 1,5$ млн маневров) ключевой проблемой является выбор алгоритмической парадигмы управления агентами. В рамках исследования был проведен сравнительный анализ применимости современных подходов:

\begin{enumerate}
    \item \textbf{Мультиагентное обучение с подкреплением (Multi-Agent Reinforcement Learning, MARL).} 
    В рамках данного подхода каждый автомобиль рассматривается как автономный интеллектуальный агент, обучающийся оптимальному выбору пути на микроуровне. Главным ограничением MARL является \textit{нестационарность среды} (non-stationarity): одновременное изменение стратегий десятками тысяч агентов делает невозможной сходимость алгоритмов обучения (например, PPO или DQN) без привлечения суперкомпьютерных мощностей. Кроме того, накладные расходы на инференс нейросетевых моделей в горячем цикле симуляции вызывают непреодолимый барьер по процессорному времени (Memory Wall).
    
    \item \textbf{Игры среднего поля (Mean Field Games, MFG).} 
    Данный аппарат сводит взаимодействие бесконечного числа агентов к системе сопряженных дифференциальных уравнений в частных производных (уравнение Гамильтона-Якоби-Беллмана и уравнение Фоккера-Планка). Несмотря на математическую строгость, MFG требует инверсии гигантских матриц Якоби для всей городской сети, что занимает минуты физического времени и исключает адаптацию к динамическим перекрытиям в реальном времени.
    
    \item \textbf{Мультиагентный поиск путей (Multi-Agent Path Finding, MAPF).}
    Классические MAPF-алгоритмы (CBS, PIBT) гарантируют бесколлизионное движение в дискретном пространстве-времени. Однако переход к Edge-based графу реального города с огромным разбросом длин ребер (от 10 м до 5 км) порождает экспоненциальный взрыв пространства состояний, делая поиск неразрешимым за приемлемое время.
\end{enumerate}

В этой связи в разрабатываемой системе обоснован и применен \textbf{гибридный агент-ориентированный подход (Agent-Oriented Approach)}, базирующийся на строгой изоляции вычислительных контуров и декомпозиции интеллекта:

\begin{itemize}
    \item \textbf{Агент полностью лишен «индивидуального интеллекта» и автономии.} Он представляет собой ультракомпактную пассивную структуру данных («квант» или «пакет» потока), движущуюся по предопределенному пути в соответствии с законами линейной 1D-кинематики.
    \item \textbf{Вычислительное ядро (Роутер \texttt{TdAltRouter})} является stateless-функцией, осуществляющей сверхбыстрый поиск путей на основе абстрактных динамических весов.
    \item \textbf{Интеллект системы вынесен на макроскопический уровень} — в Модуль принятия решений (МПР). МПР выступает в роли глобального регулятора (диспетчера), направляющего эгоистичные потоки агентов к системно-оптимальному распределению.
\end{itemize}

Такое разделение позволяет элегантно интегрировать обучение с подкреплением (RL) не на уровне отдельных автомобилей, а на **макроскопическом уровне управления средой (МПР)**. Средой (State) для RL-контроллера выступает матрица загруженности временных корзинок \texttt{volume\_buckets} дорожной сети, Действием (Action) — динамическое квотирование пропускной способности и управление директивными штрафами $mpr\_penalty$ для перегруженных артерий (Pigou congestion pricing), а Наградой (Reward) — минимизация глобального индекса задержки сети $TTI_{global}$ и предотвращение Spillback-дедлоков.

\section{Требования к динамической навигации и допущения мезоскопического моделирования}

Для обеспечения пропускной способности маршрутизатора на уровне более 5000 RPS сформулированы следующие требования к системе навигации:
\begin{enumerate}
    \item \textbf{Stateless-архитектура роутера} — полное отсутствие мьютексов и блокировок в горячем вычислительном цикле;
    \item \textbf{Zero-copy доступ к топологии} графа через системные вызовы отображения файлов в память (\texttt{mmap});
    \item \textbf{Поддержка динамического веса ребер} во временном измерении без инвалидации предрассчитанных пространственных индексов.
\end{enumerate}

Специфика мезоскопического моделирования транспортных потоков в дорожных сетях сверхбольшой размерности накладывает ряд системных допущений:
\begin{itemize}
    \item Отказ от микроскопического моделирования (декартовых координат $x, y$, угла поворота автомобиля, полосности, перестроений и безопасной дистанции следования);
    \item Пространство дискретизируется до направленных ребер графа, а время — до фиксированных 5-минутных интервалов (Volume Buckets);
    \item Движение агента описывается эйлеровым подходом: симулятор оценивает скалярную плотность потока в сечении ребра, а не траектории отдельных машин.
\end{itemize}

\section{Теоретическое обоснование мезоскопической модели транспортного потока}

\subsection{Эйлеров подход и модель пространственных очередей}
В отличие от классического лагранжева подхода (SUMO), отслеживающего координаты каждого автомобиля с вычислительной сложностью разрешения коллизий $O(N^2)$, в разработанной системе применен \textbf{эйлеров подход}. Непрерывное физическое пространство проецируется на дискретные направленные сегменты дорожной сети. Движение описывается моделью пространственных очередей (Queue-based kinematics). Состояние сегмента $e$ характеризуется его вместимостью $C_{max}$, текущим количеством находящихся на нем агентов $V(t)$ и временем свободного проезда $T_{free}$.

\subsection{Теория Spillback и гидродинамическое приближение LWR}
Для моделирования физики заторов применяется дискретная аппроксимация гидродинамической теории Lighthill-Whitham-Richards (LWR). Когда количество ТС на ребре достигает критической плотности (Jam Capacity, $V(t) \ge C_{max}$), пропускная способность входящего сечения падает до нуля, блокируя переход агентов с предыдущих ребер. Это моделирует эффект Spillback — распространение волны затора назад по топологии графа. 

Расчеты плотности затора ведутся непосредственно в **Agent Space** (без масштабирования на Agent Scale Factor), что предотвращает накопление численных дедлоков и снижает долю «черных зон» (тупиков заторов) в сети с 22,3\% до 5,7\%.

\subsection{Динамические функции сопротивления сети BPR}
Связь между плотностью потока и временной задержкой описывается модифицированной функцией Бюро дорог общего пользования США (BPR):
\begin{equation}
    W_{total}(t) = T_{free} + \Delta W_{dynamic}(t) + mpr\_penalty
\end{equation}
где $T_{free}$ — статическая нижняя граница времени проезда (включающая длину сегмента, скоростной лимит и предрассчитанный штраф за геометрический маневр), $mpr\_penalty$ — директивный штраф МПР, а $\Delta W_{dynamic}(t)$ — динамическая задержка затора:
\begin{equation}
    \Delta W_{dynamic}(t) = T_{free} \cdot \alpha \cdot \left(\frac{V(t)}{C_{300}}\right)^2
\end{equation}

Для утилизации конвейера процессора и исключения деления в горячем цикле, постоянные коэффициенты ребра сворачиваются в предрассчитанный 32-битный множитель $K_{magic}$ с применением fixed-point арифметики (битовый сдвиг на 20 разрядов влево):
\begin{equation}
    K_{magic} = \left\lfloor \frac{T_{free} \cdot \left( \frac{V_{free}}{5} - 1 \right)}{C_{300}^2} \cdot 2^{20} \right\rfloor
\end{equation}
где крутизна деградации скорости $\alpha = \frac{V_{free}}{5} - 1$ аналитически компенсирует глобально зафиксированную степень $\beta = 2$ для локальных улиц. В горячем цикле расчет BPR-штрафа сводится к 3 тактам процессора:
\begin{equation}
    \Delta W_{dynamic}(t) = (K_{magic} \cdot V(t)^2) \gg 20
\end{equation}

\section{Математический аппарат поиска путей с временной зависимостью}

\subsection{Time-Dependent A* и кольцевой буфер корзинок}
Поиск пути осуществляется на графе с динамическими весами ребер, зависящими от времени прибытия агента на сегмент. Горизонт прогнозирования в 2 часа разбит на 24 интервала по 5 минут (300 секунд). Индекс корзинки прогнозирования $t\_idx$ рассчитывается без блокировок через кольцевой буфер по абсолютному времени прибытия:
\begin{equation}
    t\_idx = \left( \left\lfloor \frac{arrival\_time}{300} \right\rfloor + head\_index \right) \bmod 24
\end{equation}
где $head\_index$ смещается асинхронным Менеджером времени каждые 5 минут с параллельной очисткой устаревшего хвоста буфера через быстрый \texttt{memset}.

\subsection{Доказательство допустимости эвристики ALT}
Эвристика A* ALT базируется на неравенстве треугольника. Для набора маяков $L$ и весов свободного потока $T_{free}$ эвристическое расстояние до цели $v$ вычисляется как:
\begin{equation}
    H_{alt}(u, v) = \max_{L} |T_{free}(L, v) - T_{free}(L, u)|
\end{equation}
Согласно неравенству треугольника, разность расстояний до маяка никогда не превышает прямое расстояние на пустом графе: $H_{alt}(u, v) \le T_{free}(u, v)$. 

Поскольку динамический штраф затора всегда строго неотрицателен ($\Delta W_{dynamic} \ge 0$), реальное время проезда по нагруженной сети всегда больше или равно времени свободного потока: $W_{real}(u, v, t) \ge T_{free}(u, v)$. Отсюда следует строгое доказательство допустимости эвристики в любой момент времени $t$:
\begin{equation}
    H_{alt}(u, v) \le T_{free}(u, v) \le W_{real}(u, v, t)
\end{equation}
Эвристика никогда не переоценивает реальную стоимость пути, гарантируя нахождение глобального математического оптимума без необходимости инвалидации матрицы маяков при возникновении заторов.

\subsection{Свойство FIFO и линейная интерполяция LERP}
Дискретизация времени на 5-минутные интервалы порождает ступенчатые функции весов, что нарушает принцип FIFO ( First-In-First-Out) — агент, выехавший позже, может обогнать выехавшего раньше за счет резкого падения штрафа в новом слоте.

Для соблюдения FIFO функция стоимости должна быть Липшиц-непрерывной со скоростью падения веса не более скорости течения времени: $\partial W_{total} / \partial t \ge -1$. Это достигается применением механизма линейной интерполяции (LERP) на границах корзинок. В горячем цикле интерполируются непосредственно вычисленные BPR-штрафы соседних слотов $P_1$ и $P_2$:
\begin{equation}
    \Delta W_{dynamic}(t) = \left( P_1 + \frac{(P_2 - P_1) \cdot (arrival\_time \bmod 300)}{300} \right) \gg 20
\end{equation}
Макроскопический закон сохранения ТС лимитирует максимальный перепад объемов между слотами, а алгоритм нарезки перекрестков препроцессинга жестко ограничивает максимальную длину ребер $T_{free}$, что математически и аппаратно гарантирует выполнение условия FIFO во всей сети.

\section{Выводы по первой главе}
Разработанная математическая модель мезоскопической навигации переводит транспортную систему от эгоистичного Пользовательского равновесия (User Equilibrium) к Системному оптимуму (System Optimum) за счет механизма упреждающего резервирования емкостей. Доказана допустимость ALT эвристики и строгое выполнение условия FIFO через LERP-интерполяцию. Обоснован отказ от тяжелых интеллектуальных агентов в пользу пассивных 8-байтовых структур данных с выносом контура обучения с подкреплением (RL) на макроскопический уровень МПР, функционирующий в сверхбыстром Fluid-режиме симуляции.
